\documentclass{dlproblemset}
\usepackage{booktabs}
\usepackage{hyperref}
\usepackage{tikz}
\usetikzlibrary{arrows.meta,decorations.pathreplacing,positioning}
\usepackage{emoji}
\tikzset{every picture/.style={line width=0.75pt}}

\definecolor{CellBlue}{HTML}{BBDEFB}
\definecolor{CellGreen}{HTML}{C8E6C9}
\definecolor{CellYellow}{HTML}{FFF9C4}
\definecolor{CellRed}{HTML}{FFCDD2}
\definecolor{LightGray}{HTML}{F5F5F5}
\definecolor{DarkBlue}{HTML}{1565C0}
\definecolor{DarkGreen}{HTML}{1B5E20}
\definecolor{DarkRed}{HTML}{C62828}

\title{Treinamento de LLMs}
\subtitle{Pré-treinamento, Alinhamento, Otimizações e PEFT}

\begin{document}

\paragraph{Instruções:} Cada questão possui exatamente uma alternativa correta.

% ============================================================
\section*{Questões}
% ============================================================

\begin{enumerate}[I)]

% ------------------------------------------------------------------
% Q1 — Fácil (Cosine Decay: cálculo numérico do LR)
% ------------------------------------------------------------------
\item Um escalonador \textit{cosine decay} sem \textit{warm-up} é configurado com
      $\eta_{\max} = 10^{-3}$, $\eta_{\min} = 10^{-4}$ e duração total
      $T = 10\,000$ passos, seguindo a fórmula vista em aula:
      \[
        \eta_t = \eta_{\min} + \frac{\eta_{\max}-\eta_{\min}}{2}\!\left(1+\cos\frac{\pi t}{T}\right).
      \]
      Qual o valor de $\eta_t$ na metade do treinamento ($t = 5\,000$)?

\begin{enumerate}[a)]
    \item $5.5 \times 10^{-4}$.
    \item $1.0 \times 10^{-3}$.
    \item $4.5 \times 10^{-4}$.
    \item $1.0 \times 10^{-4}$.
\end{enumerate}
% R: a

% ------------------------------------------------------------------
% Q2 — Fácil (memória de fine-tuning completo em mixed precision)
% ------------------------------------------------------------------
\item Um modelo possui $N = 2 \times 10^9$ parâmetros (2B). Um
      \textit{fine-tuning} completo com AdamW em \textit{mixed precision}
      mantém simultaneamente em VRAM:
      \begin{itemize}
        \item pesos em \texttt{fp16} ($2$ \texttt{B}/par)
        \item gradientes em \texttt{fp16} ($2$ \texttt{B}/par)
        \item cópia mestre dos pesos em \texttt{fp32} ($4$ \texttt{B}/par)
        \item primeiro momento Adam $\mathbf{m}$ em \texttt{fp32} ($4$ \texttt{B}/par)
        \item segundo momento Adam $\mathbf{v}$ em \texttt{fp32} ($4$ \texttt{B}/par)
      \end{itemize}
      Qual o total de VRAM requerido por esses componentes (sem contar ativações)?

\begin{enumerate}[a)]
    \item $32$ \texttt{GB}.
    \item $4$ \texttt{GB}.
    \item $8$ \texttt{GB}.
    \item $24$ \texttt{GB}.
\end{enumerate}
% R: a

% ------------------------------------------------------------------
% Q3 — Médio (SFT: papel da máscara do prompt)
% ------------------------------------------------------------------
\item A perda do SFT é
      \[
        \mathcal{L}_{\text{SFT}}(\theta) = -\sum_{t \in \mathcal{R}}
        \log P_\theta(x_t \mid x_{<t}),
      \]
      ou seja, somente os tokens da resposta ($\mathcal{R}$) contribuem. Um estudante
      sugere calcular a perda em \emph{todos} os tokens (prompt $+$ resposta),
      argumentando que ``mais supervisão $=$ melhor aprendizado''. Qual a
      principal falha desse raciocínio?

\begin{enumerate}[a)]
    \item O gradiente passa a incluir contribuições dos tokens de
          instrução, que seguem formatos repetitivos em todos os exemplos;
          o modelo gastaria capacidade aprendendo a \emph{gerar
          instruções} em vez de aprender a \emph{responder} a elas.
    \item A perda ficaria numericamente indefinida porque os tokens do prompt
          não pertencem ao vocabulário do modelo.
    \item O \textit{teacher forcing} falharia, pois tokens especiais como
          \texttt{<|im\_start|>} não podem ser previstos.
    \item A normalização por $|\mathcal{R}|$ resultaria em divisão por zero,
          já que todos os tokens passariam a contar.
\end{enumerate}
% R: a

% ------------------------------------------------------------------
% Q4 — Médio (BF16 vs. FP16: faixa dinâmica)
% ------------------------------------------------------------------
\item Uma equipe treina um modelo de 7B em \texttt{fp16} \textit{mixed precision}
      e observa \texttt{NaN} frequentes nos gradientes. Ao trocar para
      \texttt{bf16}, o problema desaparece sem precisar de \textit{loss scaling}.
      Qual a explicação correta para esse comportamento?

\begin{enumerate}[a)]
    \item \texttt{BF16} tem $8$ bits de expoente (mesma faixa dinâmica do
          \texttt{fp32}: $\sim\!10^{-38}$ a $\sim\!10^{38}$), enquanto
          \texttt{fp16} tem $5$ bits de expoente e satura próximo de
          $65\,504$. Gradientes de redes profundas extrapolam esse limite,
          causando \texttt{Inf}/\texttt{NaN}; a faixa larga do \texttt{bf16}
          elimina o \textit{overflow} sem necessidade de escalonamento.
    \item \texttt{BF16} tem $16$ bits de mantissa, oferecendo mais precisão
          que os $10$ bits do \texttt{fp16}, eliminando os erros de
          arredondamento que causavam os \texttt{NaN}.
    \item \texttt{BF16} é \textit{lossless} para pesos de redes neurais, ao
          contrário do \texttt{fp16}, que sempre introduz erros de
          quantização.
    \item \texttt{BF16} usa o dobro de bits do \texttt{fp16} ($32$ vs.\ $16$),
          o que explica o ganho de estabilidade numérica.
\end{enumerate}
% R: a

% ------------------------------------------------------------------
% Q5 — Médio (Reward Model: perda Bradley-Terry numérica)
% ------------------------------------------------------------------
\item O \textit{Reward Model} (RM) é treinado com a perda de Bradley-Terry:
      \[
        \mathcal{L}_{\text{RM}} = -\log\sigma\!\big(r_\phi(x,y_w) - r_\phi(x,y_l)\big).
      \]
      Para um par de treino, o RM atribui $r_\phi(x,y_w) = 1.8$ e
      $r_\phi(x,y_l) = 0.3$. Usando $\sigma(1.5) \approx 0.818$,
      qual o valor de $\mathcal{L}_{\text{RM}}$?

\begin{enumerate}[a)]
    \item $\approx 0.201$.
    \item $\approx 1.702$.
    \item $\approx 0.693$.
    \item $\approx 1.5$.
\end{enumerate}
% R: a

% ------------------------------------------------------------------
% Q6 — Médio (DPO vs. PPO: contabilidade de memória)
% ------------------------------------------------------------------
\item Durante o alinhamento de um modelo de 7B, o RLHF-PPO mantém
      \textbf{4 modelos} simultâneos em VRAM (política, referência SFT,
      \textit{reward model}, \textit{value function}); o DPO mantém apenas
      \textbf{2} (política $+$ referência congelada). Suponha cada cópia
      \texttt{bf16} $\approx 14$ \texttt{GB} e que o estado do otimizador
      AdamW apenas da \emph{política} ocupa $\approx 100$ \texttt{GB}
      (ignorando ativações). Quanto cada método consome e qual a economia
      do DPO?

\begin{enumerate}[a)]
    \item PPO $\approx 156$ \texttt{GB}, DPO $\approx 128$ \texttt{GB} (economia $\approx 28$ \texttt{GB}).
    \item PPO $\approx 114$ \texttt{GB}, DPO $\approx 114$ \texttt{GB} (economia $\approx 0$).
    \item PPO $\approx 56$ \texttt{GB}, DPO $\approx 28$ \texttt{GB} (economia $\approx 28$ \texttt{GB}).
    \item PPO $\approx 28$ \texttt{GB}, DPO $\approx 14$ \texttt{GB} (economia $\approx 14$ \texttt{GB}).
\end{enumerate}
% R: a

% ------------------------------------------------------------------
% Q7 — Médio (LoRA: contagem de parâmetros)
% ------------------------------------------------------------------
\item Um \textit{transformer} com $d_{\text{model}} = 2\,048$ possui quatro
      projeções de atenção por camada ($Q, K, V, O$), cada uma de formato
      $2\,048 \times 2\,048$. Aplicando LoRA com \textit{rank} $r = 16$ nas
      quatro projeções ($\mathbf{h} = \mathbf{W}_0 \mathbf{x} + \tfrac{\alpha}{r}\,\mathbf{B} \mathbf{A} \mathbf{x}$, com
      $\mathbf{B} \in \mathbb{R}^{d \times r}$, $\mathbf{A} \in \mathbb{R}^{r \times k}$),
      quantos parâmetros treináveis são adicionados por camada e que
      fração representam dos pesos originais de atenção?

\begin{enumerate}[a)]
    \item $262\,144$ params ($\approx 1.56\%$ dos $16.78\,$M originais).
    \item $131\,072$ params ($\approx 0.78\%$).
    \item $524\,288$ params ($\approx 3.12\%$).
    \item $65\,536$ params ($\approx 0.39\%$).
\end{enumerate}
% R: a

% ------------------------------------------------------------------
% Q8 — Médio (LoRA: por que $B = 0$ na inicialização?)
% ------------------------------------------------------------------
\item No LoRA, $\mathbf{A}$ é inicializado com distribuição normal gaussiana e
      $\mathbf{B}$ com \textbf{zeros}. Um estudante pergunta: ``por que não
      inicializar \emph{ambos} aleatoriamente?''. Qual a justificativa
      correta para a escolha $\mathbf{B} = \mathbf{0}$?

\begin{enumerate}[a)]
    \item Com $\mathbf{B} = \mathbf{0}$, temos $\Delta\mathbf{W} = \tfrac{\alpha}{r}\mathbf{B}\mathbf{A} = \mathbf{0}$ no passo
          inicial, logo o modelo adaptado começa \emph{exatamente} do
          pré-treinamento. Se $\mathbf{B}$ fosse aleatório, $\Delta\mathbf{W} \neq \mathbf{0}$
          adicionaria ruído às projeções pré-treinadas, degradando o ponto
          de partida.
    \item Inicializar $\mathbf{B}$ aleatoriamente faria com que o adaptador LoRA
          competisse com os pesos originais, reduzindo o número de
          parâmetros treináveis efetivos.
    \item $\mathbf{B} = \mathbf{0}$ é matematicamente necessário para que a decomposição
          $\Delta\mathbf{W} = \mathbf{B}\mathbf{A}$ tenha \textit{rank} exatamente $r$: inicialização
          aleatória quebraria a restrição de baixo \textit{rank}.
    \item Com $\mathbf{B}$ aleatório, o fator $\alpha/r$ causaria explosão de
          gradiente nas primeiras iterações.
\end{enumerate}
% R: a

% ------------------------------------------------------------------
% Q9 — Médio (Flash Attention: online softmax numérico)
% ------------------------------------------------------------------
\item Um bloco do Flash Attention processa os escores $\mathbf{s} = [1,\,2,\,0,\,3]$
      em dois tiles: $\mathbf{s}^{(1)} = [1,\,2]$ e $\mathbf{s}^{(2)} = [0,\,3]$. A recursão do
      \textit{online softmax} é
      \[
        m^{\text{new}} = \max(m^{\text{old}},\,\max_j s_j),
        \qquad
        \ell^{\text{new}} = e^{m^{\text{old}}-m^{\text{new}}}\,\ell^{\text{old}}
          + \textstyle\sum_j e^{s_j - m^{\text{new}}}.
      \]
      Qual o valor de $\ell$ após processar os dois tiles?
      (Use $e^{-1} \approx 0.368$, $e^{-2} \approx 0.135$,
      $e^{-3} \approx 0.050$.)

\begin{enumerate}[a)]
    \item $\approx 1.553$.
    \item $\approx 2.418$.
    \item $\approx 1.368$.
    \item $\approx 1.050$.
\end{enumerate}
% R: a

% ------------------------------------------------------------------
% Q10 — Médio (SFT com LR de pré-treinamento: forgetting)
% ------------------------------------------------------------------
\item Uma equipe faz SFT de um modelo usando $\eta = 3 \times 10^{-4}$
      (mesma LR usada no pré-treino) por 10 épocas. A perda SFT no conjunto
      de treino cai normalmente, mas \textit{benchmarks} gerais (MMLU,
      HellaSwag) \textbf{degradam visivelmente} ao longo das épocas.
      Qual diagnóstico e correção são adequados?

\begin{enumerate}[a)]
    \item \textit{Catastrophic forgetting}: LR alto por muitas épocas
          sobrescreve representações adquiridas no pré-treinamento, mesmo
          com a perda SFT em queda. Correção: reduzir LR para
          $10^{-6}$--$2\times 10^{-5}$ e limitar a $1$--$3$ épocas.
    \item Comportamento esperado: SFT sempre descarta conhecimento do
          pré-treinamento; a correção é treinar do zero com os dados SFT.
    \item \textit{Overfitting} nos próprios \textit{benchmarks}; a correção
          é adicionar mais \textit{benchmarks} ao conjunto de avaliação.
    \item Explosão de gradiente causada pelo LR alto; a correção é
          aumentar o \textit{weight decay} para $1.0$.
\end{enumerate}
% R: a

% ------------------------------------------------------------------
% Q11 — Difícil (GRPO: vantagem normalizada por grupo)
% ------------------------------------------------------------------
\item \emoji{skull-and-crossbones} Em uma iteração de GRPO para um problema
      com \textit{reward} binário ($1 = $ correto, $0 = $ errado),
      um prompt produz um grupo de $G = 6$ respostas com recompensas
      $\mathbf{r} = [1,\,1,\,0,\,1,\,0,\,1]$. A vantagem é
      \[
        \hat{A}_i = \frac{r_i - \overline{r}}{\sigma_r},
      \]
      onde $\overline{r}$ e $\sigma_r$ são a média e o desvio-padrão do
      grupo. Quais os valores de $\hat{A}$ para as respostas corretas e
      incorretas?

\begin{enumerate}[a)]
    \item $\hat{A}_{\text{correto}} \approx +0.707$;
          $\hat{A}_{\text{errado}} \approx -1.414$.
    \item $\hat{A}_{\text{correto}} = +1.000$;
          $\hat{A}_{\text{errado}} = -1.000$.
    \item $\hat{A}_{\text{correto}} = +1/3$;
          $\hat{A}_{\text{errado}} = -2/3$.
    \item $\hat{A}_{\text{correto}} = +0.500$;
          $\hat{A}_{\text{errado}} = 0$.
\end{enumerate}
% R: a

% ------------------------------------------------------------------
% Q12 — Difícil (DPO: cálculo numérico da perda)
% ------------------------------------------------------------------
\item \emoji{skull-and-crossbones} A perda do DPO é
      \[
        \mathcal{L}_{\text{DPO}} = -\log\sigma\!\left(
          \beta\log\frac{\pi_\theta(y_w|x)}{\pi_{\text{ref}}(y_w|x)}
          - \beta\log\frac{\pi_\theta(y_l|x)}{\pi_{\text{ref}}(y_l|x)}
        \right).
      \]
      Em um exemplo com $\beta = 0.1$,
      $\log\tfrac{\pi_\theta(y_w|x)}{\pi_{\text{ref}}(y_w|x)} = 2$ e
      $\log\tfrac{\pi_\theta(y_l|x)}{\pi_{\text{ref}}(y_l|x)} = -1$, qual o
      valor de $\mathcal{L}_{\text{DPO}}$?
      (Use $\sigma(0.3) \approx 0.574$, $\sigma(3) \approx 0.953$,
      $\sigma(-0.3) \approx 0.426$.)

\begin{enumerate}[a)]
    \item $\approx 0.555$.
    \item $\approx 0.048$.
    \item $\approx 0.853$.
    \item $\approx 0.300$.
\end{enumerate}
% R: a

\end{enumerate}


% ============================================================
\section*{Soluções}
% ============================================================

% --- Solução Q1 ---
\subsection*{I — Resposta: (a)}
Em $t = T/2$: $\cos(\pi \cdot 0.5) = \cos(\pi/2) = 0$, logo
$1 + \cos(\pi/2) = 1$. Portanto,
\[
  \eta_{T/2} = 10^{-4} + \tfrac{10^{-3} - 10^{-4}}{2}\cdot 1
             = 10^{-4} + 4.5\times 10^{-4}
             = 5.5\times 10^{-4}.
\]
A opção~(b) ignora totalmente o termo com cosseno.
A opção~(c) omite a parcela $\eta_{\min}$ da fórmula.
A opção~(d) confunde $\cos(\pi/2) = 0$ com $\cos(\pi) = -1$, obtendo
$\eta_{\min}$, valor correto apenas no fim do treino ($t = T$).

% --- Solução Q2 ---
\subsection*{II — Resposta: (a)}
Somando os componentes em bytes por parâmetro:
$2 + 2 + 4 + 4 + 4 = 16$ \texttt{B}/par.
Aplicando a $N = 2\times 10^9$:
$2\times 10^9 \times 16 = 3.2\times 10^{10}$ bytes $= 32$ \texttt{GB}.
As opções~(b), (c) e (d) ignoram subconjuntos dos componentes: somente pesos,
pesos+gradientes e um dos momentos do Adam, respectivamente.

% --- Solução Q3 ---
\subsection*{III — Resposta: (a)}
Os tokens de instrução seguem padrões repetitivos (\texttt{<|im\_start|>user},
``Explique...'', etc.). Se forem incluídos na perda, o gradiente é dominado
por eles e o modelo gasta capacidade modelando a \emph{distribuição das
instruções} em vez de aprender a gerar boas respostas. Mascarar o prompt
garante que o aprendizado se concentre em $P(\text{resposta} \mid \text{instrução})$,
não em $P(\text{instrução})$.
A opção~(b) é falsa: tokens de prompt estão no vocabulário (caso contrário,
não poderiam sequer ser alimentados como entrada).
A opção~(c) confunde mascaramento de \emph{loss} com ausência do token no
vocabulário.
A opção~(d) é falsa: a normalização é por $|\mathcal{R}|$ (tokens da resposta),
e incluir os tokens do prompt apenas aumentaria $|\mathcal{R}|$, não o zeraria.

% --- Solução Q4 ---
\subsection*{IV — Resposta: (a)}
A diferença essencial entre \texttt{fp16} e \texttt{bf16} é o número de bits
de expoente: \texttt{fp16} usa $5$ (limite $\approx 65\,504$) e \texttt{bf16}
usa $8$ (mesmo limite do \texttt{fp32}, $\approx 3.4\times 10^{38}$).
Gradientes grandes provocam \textit{overflow} em \texttt{fp16}, gerando
\texttt{Inf}/\texttt{NaN}; a faixa larga do \texttt{bf16} elimina esse
problema sem exigir \textit{loss scaling}.
A opção~(b) inverte o trade-off: \texttt{bf16} tem \emph{menos} mantissa
($7$ bits) que o \texttt{fp16} ($10$). O ganho não é precisão, é faixa.
A opção~(c) é falsa: ambos são lossy.
A opção~(d) é falsa: ambos ocupam $16$ bits.

% --- Solução Q5 ---
\subsection*{V — Resposta: (a)}
$\Delta r = r_\phi(x,y_w) - r_\phi(x,y_l) = 1.8 - 0.3 = 1.5$.
$\sigma(1.5) \approx 0.818$, portanto
$\mathcal{L}_{\text{RM}} = -\ln(0.818) \approx 0.201$.
A opção~(b) inverte o sinal da diferença, aplicando $\sigma(-1.5) \approx 0.182$
($-\ln(0.182) \approx 1.702$): seria a perda se o RM ordenasse
\emph{incorretamente}.
A opção~(c) usa $\sigma(0) = 0.5$, equivalente a ignorar totalmente
$\Delta r$.
A opção~(d) retorna o \textit{logit} sem passá-lo pelo $\sigma$ nem aplicar
$-\ln$.

% --- Solução Q6 ---
\subsection*{VI — Resposta: (a)}
\textbf{PPO:} 4 cópias em \texttt{bf16} $+$ estado do otimizador da política
$= 4\cdot 14 + 100 = 56 + 100 = 156$ \texttt{GB}.\\
\textbf{DPO:} 2 cópias em \texttt{bf16} $+$ estado do otimizador da política
$= 2\cdot 14 + 100 = 28 + 100 = 128$ \texttt{GB}.\\
\textbf{Economia:} $156 - 128 = 28$ \texttt{GB}, exatamente o custo das duas
cópias que o DPO dispensa (\textit{reward model} e \textit{value function}).
A opção~(b) é falsa: o PPO \emph{inclui} duas cópias adicionais frozen,
portanto sua memória é maior.
A opção~(c) ignora o estado do otimizador da política ($100$ \texttt{GB}).
A opção~(d) reduz incorretamente cada método a uma única cópia.

% --- Solução Q7 ---
\subsection*{VII — Resposta: (a)}
Cada projeção LoRA adiciona $|\mathbf{B}| + |\mathbf{A}| = d\cdot r + r\cdot k = 2\cdot d\cdot r$
parâmetros (como $d = k = 2\,048$). Para quatro projeções:
\[
  4 \times 2 \cdot 2\,048 \cdot 16 = 4 \times 65\,536 = 262\,144 \;\text{params}.
\]
Pesos originais: $4 \times 2\,048^2 = 16\,777\,216 \approx 16.78\,$M.
Fração: $262\,144 / 16\,777\,216 \approx 1.56\%$.
A opção~(b) conta apenas $d \cdot r$ por projeção, esquecendo a segunda matriz.
A opção~(c) duplica erroneamente o resultado.
A opção~(d) considera apenas uma das quatro projeções.

% --- Solução Q8 ---
\subsection*{VIII — Resposta: (a)}
O ponto de partida do LoRA é o modelo pré-treinado intacto: exige-se
$\Delta\mathbf{W} = \tfrac{\alpha}{r}\mathbf{B}\mathbf{A} = \mathbf{0}$ na inicialização. Zerar $\mathbf{B}$ (com $\mathbf{A}$
aleatório) satisfaz essa condição. Se \emph{ambos} fossem aleatórios,
$\Delta\mathbf{W} \neq \mathbf{0}$ adicionaria ruído aos pesos congelados, destruindo o
ponto de partida calibrado do pré-treinamento. O modelo teria de
``reaprender'' antes mesmo de começar a se adaptar aos dados de ajuste.
A opção~(b) não faz sentido dimensional: inicialização não altera contagem
de parâmetros treináveis.
A opção~(c) é falsa: o \textit{rank} de $\mathbf{B}\mathbf{A}$ é determinado pela forma das
matrizes, não pelos valores iniciais.
A opção~(d) confunde o papel de $\alpha/r$ (fator de escala constante, não
amplificador de gradiente).

% --- Solução Q9 ---
\subsection*{IX — Resposta: (a)}
\textbf{Processando $\mathbf{s}^{(1)} = [1,\,2]$:}
\begin{align*}
  m^{(1)} &= \max(1,\,2) = 2 \\
  \ell^{(1)} &= e^{1-2} + e^{2-2} = 0.368 + 1 = 1.368.
\end{align*}
\textbf{Processando $\mathbf{s}^{(2)} = [0,\,3]$:}
\begin{align*}
  m^{(2)} &= \max(2,\,3) = 3 \\
  \ell^{(2)} &= e^{m^{(1)}-m^{(2)}}\cdot \ell^{(1)} + e^{0-3} + e^{3-3} \\
             &= e^{-1}\cdot 1.368 + e^{-3} + e^{0} \\
             &\approx 0.368 \cdot 1.368 + 0.050 + 1 \\
             &\approx 0.503 + 0.050 + 1 = \mathbf{1.553}.
\end{align*}
Conferência (softmax global): $e^{1-3}+e^{2-3}+e^{0-3}+e^{3-3}
= 0.135 + 0.368 + 0.050 + 1 = 1.553$. \checkmark
A opção~(b) omite o fator de rescala $e^{m^{\text{old}}-m^{\text{new}}}$,
tratando $\ell^{(1)}$ como se já estivesse normalizado para $m^{(2)}$.
A opção~(c) ignora $\mathbf{s}^{(2)}$.
A opção~(d) ignora o histórico $\ell^{(1)}$.

% --- Solução Q10 ---
\subsection*{X — Resposta: (a)}
O padrão é clássico de \textit{catastrophic forgetting}: a perda SFT
decresce porque o modelo aprende o \emph{formato} dos dados de ajuste,
mas as muitas atualizações com LR alto modificam significativamente os
pesos, apagando representações úteis para tarefas gerais. O remédio
canônico é LR muito menor que no pré-treinamento ($10^{-6}$–$2\times 10^{-5}$)
e apenas $1$–$3$ épocas.
A opção~(b) contradiz a hipótese do alinhamento superficial (LIMA):
o SFT ensina formato e estilo, não substitui o pré-treinamento.
A opção~(c) é implausível: \textit{benchmarks} não estão no conjunto
de treino, então \textit{overfitting} neles não faz sentido.
A opção~(d) descreve um sintoma diferente: explosão de gradiente geraria
\texttt{NaN} agudos, não degradação gradual. Além disso, aumentar
\textit{weight decay} não corrige explosões e pioraria o
\textit{forgetting}.

% --- Solução Q11 ---
\subsection*{XI — Resposta: (a)}
Com rewards $[1,1,0,1,0,1]$ e $G = 6$:
\begin{align*}
  \overline{r} &= \tfrac{4}{6} = \tfrac{2}{3} \approx 0.667, \\
  \sigma_r^2 &= \mathbb{E}[r^2] - \overline{r}^{\,2} = \tfrac{4}{6} - \tfrac{4}{9}
             = \tfrac{6}{9} - \tfrac{4}{9} = \tfrac{2}{9}, \\
  \sigma_r   &= \tfrac{\sqrt{2}}{3} \approx 0.471.
\end{align*}
Portanto,
\[
  \hat{A}_{\text{correto}} = \frac{1 - 2/3}{\sqrt{2}/3} = \frac{1/3}{\sqrt{2}/3}
    = \frac{1}{\sqrt{2}} \approx +0.707,
\]
\[
  \hat{A}_{\text{errado}} = \frac{0 - 2/3}{\sqrt{2}/3}
    = -\frac{2}{\sqrt{2}} = -\sqrt{2} \approx -1.414.
\]
A magnitude maior para as respostas erradas (minoritárias) é esperada: o
z-score pesa mais a cauda menos populosa.
A opção~(b) usa os valores do slide, que valem apenas quando \emph{metade}
do grupo acerta ($G = 4$, 2 corretos, 2 errados), situação em que
$\sigma_r = 0.5$.
A opção~(c) esquece a normalização por $\sigma_r$.
A opção~(d) substitui a \textit{baseline} correta (média) por $\max(r) = 1$.

% --- Solução Q12 ---
\subsection*{XII — Resposta: (a)}
Argumento do $\sigma$:
\[
  \beta\big(\log\tfrac{\pi_\theta(y_w|x)}{\pi_{\text{ref}}(y_w|x)}
         - \log\tfrac{\pi_\theta(y_l|x)}{\pi_{\text{ref}}(y_l|x)}\big)
  = 0.1 \cdot (2 - (-1)) = 0.3.
\]
$\sigma(0.3) \approx 0.574$, portanto
$\mathcal{L}_{\text{DPO}} = -\ln(0.574) \approx 0.555$.
A opção~(b) esquece o $\beta$, aplicando $\sigma(3) \approx 0.953$,
subestimando dramaticamente a perda.
A opção~(c) inverte o sinal, usando $\sigma(-0.3)$ (corresponderia a uma
\emph{preferência errada}: $y_l$ mais provável que $y_w$).
A opção~(d) retorna o argumento do $\sigma$ sem aplicar $-\log\sigma$.

\end{document}
